```latex
\documentclass[12pt]{article}
\usepackage[utf8]{inputenc}
\usepackage{amsmath, amssymb, graphicx, hyperref}
\usepackage{geometry}
\geometry{a4paper, margin=1in}
\title{SHOA-Chrono: A Unified Spintronic Framework for Quantum Memory Readout and the Chronoscope T-G Protocol}
\author{Konstantin Fedotov}
\date{\today}

\begin{document}

\maketitle

\begin{abstract}
We present SHOA-Chrono, a theoretical and engineering foundation for non-invasive readout of quantum memory in crystalline media. Building upon the SHOA-Neuro model, we formalize the Chronoscope T-G protocol, which integrates four components: (1) an event hypothesis, (2) a probing GRAPE-optimized impulse, (3) a library of archetypal patterns (ME), and (4) Bayesian reconstruction with a Laplacian sparse prior. We provide a complete experimental configuration and demonstrate a simulated test protocol on a Wardenclyffe Tower sample, reconstructing the event with 89\% accuracy.
\end{abstract}

\section{Introduction}
The SHOA project has evolved from self-healing through learning to adaptive self-assembly. SHOA-Chrono represents the next logical step: quantum memory readout. We assert that a crystalline lattice can retain information about past events as stable spin correlations, and this information can be extracted via resonant interaction with a probing impulse.

\section{Mathematical Framework: The T-G Protocol}
The Chronoscope T-G protocol is formalized as a controlled evolution of the density matrix for an ensemble of crystal defects under a GRAPE-optimized impulse. The total Hamiltonian is:
\begin{equation}
\mathcal{H}_{\text{Chrono}} = \mathcal{H}_{\text{defect}} + \mathcal{H}_{\text{dip}} + \mathcal{H}_{\text{probe}}(t) + \mathcal{H}_{\text{env}},
\end{equation}
where $\mathcal{H}_{\text{probe}}(t)$ is modulated by the Grape-cluster phase gate.

The readout operator linking neural activity with SHOA-resilience is:
\begin{equation}
\hat{J}_{\text{readout}}(t) = \frac{d}{dt} \sum_k w_k C_k(t),
\end{equation}
where $w_k$ are weights proportional to the SHOA-resilience of the $k$-th domain and $C_k(t)$ are measured cross-correlation functions.

\section{Experimental Configuration and Simulation}
We fully specify the experimental setup: SNSPD detector, 4.0 K, Sumerian\_Base\_3.0 ME library, quartz parameters. A simulated scan of a Wardenclyffe sample reconstructed a humanoid figure with raised arm at confidence 0.91, consistent with the Tesla full-power test hypothesis.

\section{Conclusion}
SHOA-Chrono and the T-G protocol represent a platform fully ready for verification. We hand it over to the global scientific community, calling for independent reproduction of the results.

\end{document}
